\section{Conclusions}
\label{sec:concl}

% What was achieved?
This work set out to operate a USV and a UAV connected by a tether as a single system at sea, each vehicle within its own limits and both within the reach of the tether. Models of both vehicles, derived from their physics and validated against the simulator, predict their motion over the horizons of the controllers with errors of a few percent, and a closed-form model of a slack tether that accounts for its weight reduces the error of its force by about three quarters relative to a taut tether or a hanging load. On these models, each vehicle is driven by its own NMPC, and the two exchange the trajectories they predict, so that each keeps the separation within the reach of the tether while anticipating where the other vehicle is going.

% What do the results show?
Both controllers track the trajectories given to their vehicles within a few centimetres, the UAV more than four times closer than the position controller of PX4, and with nearly the same error under a wind of $8\,\text{m/s}$, which the disturbance observer compensates. When a reference demands more than a vehicle can deliver, each controller gives up following it only as far as the limits require, and the separation between the vehicles is held at the reach of the tether, so that any trajectory can be given to either vehicle without putting the tether at risk, and the manoeuvres can remain simple. The UAV lands on the moving deck a few centimetres from the centre of the pad in every sea tested, with a speed of contact that grows with the waves and in the higher sea reaches the default landing speed of PX4, and, with the disturbance observer, also under wind, which otherwise forces the landing to be aborted. Coupled by the tether, the two vehicles complete an inspection, a patrol and a recovery as a single system, and in the inspection the tether becomes not a constraint to be avoided, but a means to survey a tower that neither vehicle could survey alone.

% What comes next?
The system was developed and evaluated in simulation, and bringing it to real operations at sea is a long path. Its main steps are the hardware of the tethered system, with RTK receivers on both vehicles that also give their relative position, the validation of the models and the estimation of the state on the real vehicles, the handling of failures and the avoidance of obstacles, which may require longer horizons, and a gradual validation from land to the open sea, where a prediction of the motion of the deck could time the touchdown. Even then, a much larger volume of testing would be needed before real use.
